% Deutsche Parallelfassung zu introduction.tex

\chapter{Einleitung}

\section{Das Eisenbahnproblem vor dem Informationsproblem}

Ein Eisenbahnfahrzeug kann keinen benachbarten Weg wählen, wenn die Geometrie
ungünstig wird. Seine Radsätze sind gezwungen, dem Schienenpaar zu folgen. Ein
Übergang zwischen Gerade und Kreisbogen verändert daher mehr als einen
gezeichneten Verlauf: Krümmungsgesetz und Überhöhungsentwicklung bestimmen
zusammen mit der Geschwindigkeit, wie sich Querbeschleunigung, Rollbewegung,
Kraftentwicklung, Komfort, Verschleiß und Zulässigkeit entlang der
Fahrzeugtrajektorie entwickeln.

Man vergleiche eine Klothoide mit einem Übergang, bei dem auch die
Krümmungsableitung an den Enden glatt ausläuft. Beide können dieselben
Randkrümmungen verbinden. Die Klothoide liefert ein konstantes
\(\kappa'(s)\), während ein geglättetes Gesetz die Endwertsprünge von
\(\kappa'(s)\) unterdrücken und \(\kappa''(s)\) regeln kann. Letzteres kann
günstigeres physikalisches Verhalten bieten, verwendet dafür aber gewöhnlich
Länge. Im ursprünglichen Ingenieurvergleich wurde der entscheidende Vorteil
der Klothoide deshalb nicht in einer universellen dynamischen Überlegenheit
gesehen, sondern häufig in ihrer kompakten, transparenten Konstruktion. Diese
Beobachtung eröffnete die Frage, die AIM vorausgeht: Wie lassen sich
unterschiedliche bekannte Übergangsformen in einer gemeinsamen funktionalen
Sprache ausdrücken, berechnen und vergleichen?

Die ursprüngliche Antwort war \emph{Berlinish}, eine verallgemeinerte
Übergangsgrammatik, und \emph{axtranNew}, der vorgesehene Berechnungskern, der
diese Grammatik anwendbar macht. Die weitere ufAIM-Umgebung entwickelte sich,
weil Formeln und Solver-Ergebnisse allein das entstehende Wissen nicht
dauerhaft, untersuchbar, editierbar, zurechenbar und über
Ingenieurarbeitsabläufe hinweg nutzbar halten konnten. Diese Geschichte ist
eine Interpretation der Entwicklungslinie durch die Thesis, keine Behauptung,
Berlinish oder axtranNew sei ein genehmigter Kernel-Begriff.

\section{Eine Änderung, die nicht in einem Modell verbleibt}

Betrachten wir ein Alignment, das in lokalen Ingenieurkoordinaten übernommen
wurde. Ein Ingenieur ändert die Krümmung eines Übergangs. Die geometrische
Wirkung tritt unmittelbar ein: Der nachfolgende Verlauf verschiebt sich,
abgeleitete Positionen ändern sich und auf dem Bildschirm erscheint eine neue
Kurve. Doch diese scheinbar lokale Bearbeitung wirft umfassendere Fragen auf.
Beschreibt das bearbeitete Modell noch dasselbe Alignment? Welche
Beobachtungen beziehen sich auf das bestehende Gleis und welche Werte drücken
den vorgeschlagenen Entwurf aus? Welche Regeln gelten in diesem
Projektzusammenhang? Wer ist berechtigt, die Folge anzunehmen?

Keine einzelne Sicht beantwortet all diese Fragen. GIS kann das Alignment
zwischen Gelände, Flurstücken, Netzen und Umweltdaten verorten, doch aus der
Fülle des verfügbaren Kontexts muss das Relevante ausgewählt werden. CAD gibt
dem Ingenieur genaue Konstruktions- und Bearbeitungsmöglichkeiten, zeigt
üblicherweise aber nur den für die unmittelbare Aufgabe erforderlichen
Ausschnitt. BIM- und Austauschmodelle können Objekte und Attribute
fachübergreifend verbinden, doch Lieferanforderungen bestimmen häufig, was in
welcher Tiefe repräsentiert wird. Messungen, Berechnungsblätter, Normen,
Vorgangsakten und erfahrene Ingenieure tragen weitere Teile der Bedeutung.

Diese Fragmentierung ist kein Beleg für ein Versagen der Werkzeuge. Jedes ist
innerhalb seines Zwecks nützlich. Die Schwierigkeit entsteht zwischen ihren
Zwecken: Ingenieurwissen ist über Repräsentationen, Dokumente, Abläufe und
Menschen verteilt, während die Folge einer Alignment-Änderung in einem
bestimmten Handlungsmoment verstanden werden muss.

\begin{quote}
Ingenieursoftware kann nahezu jede sichtbare Folge einer Alignment-Änderung
beschreiben; häufig kann sie jedoch nicht angeben, was gleich bleiben muss,
während sich diese Beschreibungen ändern.
\end{quote}

Diese Spannung motiviert die vorliegende Thesis.

\section{Das Problem ist nicht allein der Geometrieaustausch}

Geometrische Interoperabilität ist wichtig. Die erfolgreiche Übertragung von
Koordinaten oder Kurvenelementen löst das Ingenieurproblem jedoch nicht. Sechs
Unterscheidungen müssen die Übertragung überstehen.

\begin{description}[style=nextline]
\item[Identität]
Das Engineering Object muss über Dateien, Sichten, Koordinatenrealisierungen
und Änderungen im Arbeitsablauf hinweg unterscheidbar bleiben. Numerische
Gleichheit allein belegt nicht, dass zwei Beschreibungen dasselbe Objekt
betreffen; ebenso entscheidet eine geometrische Änderung allein nicht, dass
aus einem Objekt ein anderes geworden ist.

\item[Konstruktive Absicht]
Die geordnete longitudinale Konstruktion und ihr Krümmungsverlauf erklären,
wie das Alignment definiert ist. Abgetastete Punkte oder gerenderte Kurven
können diese Konstruktion sichtbar machen, besitzen sie aber nicht.

\item[Realität und Beobachtung]
Das beabsichtigte Alignment, das physisch realisierte Gleis und eine
Beobachtung dieses Gleises sind miteinander verbunden, aber verschieden. Eine
Messung ist Evidenz, die unter bestimmten Bedingungen gewonnen wurde; sie ist
nicht automatisch der physische Zustand oder angenommenes Ingenieurwissen.

\item[Änderung und Folge]
Eine Krümmungsänderung wirkt sich auf Position, Orientierung,
Überhöhungszusammenhang, Fahrzeugreaktion, Lichträume, Bauausführung und
Regelkonformität aus. Die Abhängigkeiten, die diese Folgen übertragen, müssen
hinreichend ausdrücklich sein, um eine Bewertung zu ermöglichen.

\item[Knowledge Ownership]
Ein Format, Viewer, Solver oder eine Datenbank kann Ingenieurwissen tragen und
verwenden, ohne Autorität über dessen Bedeutung zu erwerben. Geltungsbereich,
Provenienz, Anwendbarkeit und verantwortliches Ownership müssen sichtbar
bleiben.

\item[Entscheidung]
Berechnung und Optimierung können Evidenz und Alternativen erzeugen. Sie nehmen
für sich genommen keine Ingenieuränderung an. Bewertung und
Entscheidungsautorität müssen getrennt bleiben.
\end{description}

Das ungelöste Problem lässt sich daher knapp formulieren: Wie kann ein
Alignment über Repräsentationen und Änderungen hinweg eine kohärente
ingenieurtechnische Bedeutung bewahren und zugleich das richtige Wissen
anwendbar machen, wenn ein Ingenieur handelt?

\section{Forschungsfrage und zentrale These}

Die leitende Forschungsfrage lautet:

\begin{quote}
Wie lassen sich Alignment-spezifische ingenieurtechnische Identität,
Konstruktion, Zustand, Realität, Beobachtung, Wissen und Bewertung so ordnen,
dass Änderungen nachvollziehbar bleiben und ihre Folgen über Repräsentationen
hinweg berechenbar werden?
\end{quote}

Diese Frage verallgemeinert das ursprüngliche Übergangsproblem. Die erste
technische Frage lautete, wie Krümmungsfunktionen, Randbedingungen,
Komponentenlängen und Freiheitsgrade Eisenbahnübergänge beschreiben und lösen
können, ohne eine Kurvenfamilie für allgemein überlegen zu erklären. Sobald
diese berechneten Übergänge Bearbeitung, Austausch, Beobachtung, Vergleich und
verantwortliche Verwendung überstehen mussten, wurde die
Informationsmodellierungsfrage unvermeidlich.

Die Thesis antwortet mit Alignment-Based Information Modelling (AIM). AIM ist
wissensbasiert: Die ingenieurtechnische Nutzung hängt von anerkannter Bedeutung
und Anwendbarkeit ab. AIM ist Alignment-spezifisch: Longitudinale Konstruktion
und Krümmungsverlauf strukturieren das Objekt. AIM ist
repräsentationsbewusst: Keine Datei, Sicht oder Koordinatenrealisierung wird
mit dem Objekt selbst verwechselt. AIM ist implementierungsfähig: Die
Beziehungen werden zu Zuständen, Operatoren, Nebenbedingungen, Pipelines und
numerischen Verfahren entwickelt.

\begin{proposition}[Zentrale These]
Ein Alignment kann über Entwurf, Realisierung, Beobachtung, Bewertung und
Änderung hinweg kohärent behandelt werden, wenn seine intrinsische konstruktive
Identität von seinen Repräsentationen getrennt bleibt und durch ausdrückliche
Zustands- und Operatorbeziehungen mit anwendbarem Ingenieurwissen verbunden
wird.
\end{proposition}

Krümmung besitzt in dieser Argumentation eine besondere Rolle. Bei einem
Alignment gehört der geordnete Krümmungsverlauf zur intrinsischen
konstruktiven Struktur, während Position und Pose durch metrische
Interpretation und Integration entstehen. Deshalb führt die Thesis von der
vertrauten Kurve zu einem krümmungsgetriebenen generativen Modell. Sie
verneint Geometrie nicht; sie erklärt, woher Geometrie stammt und was
Geometrie allein nicht besitzen kann.

\section{Beiträge}

Die Arbeit leistet fünf Arten von Beiträgen.

\begin{description}[style=nextline]
\item[Begrifflicher Beitrag]
Sie entwickelt eine kohärente, am Alignment ausgerichtete Darstellung von
Engineering Object Identity, konstruktiver Identität, Repräsentation,
metrischer und physischer Realisierung, Beobachtung, Ingenieurwissen, Bewertung
und Entscheidung.

\item[Formaler und methodischer Beitrag]
Sie entwickelt krümmungsgetriebene Pose- und Zustandsformulierungen,
ausdrückliche räumliche und Laufzeitoperatoren,
realisierungsbewusste Abhängigkeiten, die Bewertung von Nebenbedingungen und
strukturierte Optimierungsformulierungen.

Sie macht außerdem die ursprüngliche Berlinish-Komposition als mathematisches
Framework ausdrücklich, das Halbwellen-, Klothoiden-, asymmetrische und
zusammengesetzte Übergangsformen zueinander in Beziehung setzt, ohne diesen
Forschungsbeitrag zur Kernel-Autorität zu erheben.

\item[Ingenieurtechnischer Beitrag]
Sie verbindet Alignment-Übergänge, Überhöhung, fahrzeugbezogene Größen,
Stationierung, Koordinatenbezug, Beobachtungen, Bauauswirkungen und
Eisenbahnregeln, ohne ihre unterschiedlichen Bedeutungen in einem einzigen
Geometriemodell zusammenfallen zu lassen.

\item[Implementierungs- und Referenzbeitrag]
Sie spezifiziert Berechnungspipelines und dünn besetzte oder hybride
Repräsentationen, die Rekonstruktion, Bewertung und Optimierung unterstützen
und zugleich die Grenze zwischen kanonischer ingenieurtechnischer Bedeutung
und abgeleiteten Produkten bewahren.

Innerhalb dieses Beitrags bezeichnet axtranNew die Berechnungslinie vom Lösen
von Anschluss- und Randbedingungen bis zu editierbaren Alignment-Strukturen.
Die aktuelle Referenzimplementierung unterstützt registry-basierte
Übergangsauswertung, dünn besetzte Rekonstruktion, ausgewählte native
Bearbeitungen, Vergleiche und begrenzte Optimierungsprobleme; ein allgemeiner
produktionsreifer Optimierer bleibt künftige Arbeit.

\item[Erklärungsbeitrag]
Sie bietet eine nach Abhängigkeiten geordnete Darstellung, über die ein
technisch qualifizierter Leser von einer gewöhnlichen Alignment-Änderung zu
der begrifflichen und rechnerischen Maschinerie gelangt, die zur Bewertung
ihrer Folgen erforderlich ist.
\end{description}

Dies sind Grundlagenbeiträge. Die Thesis demonstriert Formulierungen,
Beziehungen und Implementierungswege; sie beansprucht weder, jedes
Fachwerkzeug zu ersetzen, noch ein vollständiges
Mehrkörper-Eisenbahnmodell bereitzustellen oder ingenieurtechnische
Entscheidungsautorität an Berechnung zu übertragen.

\section{Methodische Ausrichtung}

Die Untersuchung wechselt zwischen begrifflicher Trennung und konstruktiver
Demonstration. Zunächst fragt sie, um welche Art von Gegenstand es sich jeweils
handelt---Objekt, Repräsentation, Zustand, Beobachtung, Wissen oder
Entscheidung---und welche Abhängigkeiten zulässig sind. Anschließend entwickelt
sie mathematische Modelle und Operatoren, die ausgewählte Abhängigkeiten
berechenbar machen. Eisenbahnbeispiele, normbezogene Nebenbedingungen,
Datenrekonstruktion und Optimierung prüfen, ob die Unterscheidungen unter
ingenieurtechnischem Druck nützlich bleiben.

Die Argumentation schreitet von Beispielen zur Formalisierung fort.
Gleichungen werden eingeführt, wenn sie eine Abhängigkeit offenlegen oder eine
Bewertung ermöglichen, nicht als Ersatz für die Ingenieurfrage. Numerische
Lösungen werden als Kandidaten oder Evidenz innerhalb eines verantworteten
Prozesses behandelt, nicht als Antworten, die sich selbst autorisieren.

\section{Leseweg durch die Thesis}

Die Teile sind nach Fragen geordnet, die der Leser weitertragen kann.

\begin{enumerate}[leftmargin=*]
\item \emph{Warum erfordert ein Alignment mehr als eine Kurve?} Der erste Teil
  schärft das praktische Problem und die AIM-These.
\item \emph{Was muss getrennt bleiben, bevor Mathematik helfen kann?}
  Mathematische Grundlagen und Ingenieurbegriffe bestimmen die Grenzen
  zwischen Identität, Wissen, Repräsentation, Realisierung, Beobachtung und
  Entscheidung.
\item \emph{Wie erhält die intrinsische Konstruktion geometrische und
  betriebliche Bedeutung?} Geometrie, Eisenbahnzustand und dynamischer
  Eisenbahnzustand leiten Pose, Überhöhung, fahrzeugbezogene Größen und
  Zulässigkeit aus dem longitudinalen Zustand ab.
\item \emph{Wie werden Folgen während der Nutzung eines Systems bewertet?}
  Laufzeit und Bewertung machen die Anwendung von Operatoren und Wissen
  ausdrücklich.
\item \emph{Wie bleiben Entwurf, physisches Gleis und Evidenz darüber
  miteinander verbunden, ohne identisch zu werden?} Realität und Daten
  behandeln Koordinaten, Messung, Bauausführung und realisierten Zustand.
\item \emph{Wie lässt sich die Darstellung implementieren und gezielt
  verändern?} Alignment-Modellierung, System und Pipeline sowie Optimierung
  entwickeln dünn besetzte Modelle, rechnerische Abhängigkeiten und
  Lösungsverfahren. Sie zeigen außerdem, wie Berlinish, axtranNew,
  transitionDB, TransEd, SPOT und alignmentOS unterschiedliche Rollen in der
  Entwicklungsantwort einnehmen.
\item \emph{Was besteht die Begegnung mit Regeln, Austauschsystemen und
  ingenieurtechnischer Nutzung?} Ingenieurwesen und Normen, Anwendungen,
  Diskussion und Ausblick prüfen die Reichweite des Frameworks und benennen
  seine Grenzen.
\end{enumerate}

Die zu Beginn eingeführte Krümmungsänderung wird als Test wiederkehren: Was
bleibt gleich, was ändert sich, was wird beobachtet, was wird abgeleitet und
wer darf entscheiden? Der nächste Teil kehrt zunächst zu dieser Änderung
zurück, bevor eine formale Ontologie erforderlich ist.
